\chapter*{Bản Đồ Đọc Sách}
\addcontentsline{toc}{chapter}{Bản Đồ Đọc Sách}
\markboth{Bản Đồ Đọc Sách}{Bản Đồ Đọc Sách}

\begin{truyen}{Đọc Quyển XXXIII Theo Đường Nào?}{How To Travel Through This Book}
\chuhoa{N}{ếu đọc liền một mạch}, người đọc có thể thấy quyển này giống một chuỗi chuyện đối đáp châm biếm. Nhưng nếu đọc theo trục chủ đề, nó hiện ra như một bức bản đồ tâm lý học đường. Mỗi phần không chỉ nói về một nhóm người; nó nói về một kiểu xung đột. Phần 1 là xung đột giữa cảm xúc và trách nhiệm. Phần 2 là xung đột giữa kỷ luật và uy tín. Phần 3 là xung đột giữa yêu thương và kiểm soát. Phần 4 là xung đột giữa giấc mơ và năng lực thật.
\textit{(If read straight through, this book may look like a chain of satirical school dialogues. But if it is read by theme, it becomes a psychological map of school life. Each part is not only about a group of people; it is about a type of conflict. Part 1 is the conflict between emotion and responsibility. Part 2 is between discipline and credibility. Part 3 is between love and control. Part 4 is between dreams and real capacity.)}

Vì vậy, người đọc có thể chọn đường đi phù hợp với điều mình đang cần nhất. Người học sinh đang mệt có thể bắt đầu từ Phần 1. Giáo viên đang muốn chỉnh lại cách giữ lớp có thể vào Phần 2. Phụ huynh đang căng với con nên đi thẳng vào Phần 3. Học sinh cuối cấp hoặc sinh viên năm đầu sẽ thấy Phần 4 đánh trúng nhiều ảo tưởng của mình nhất.
\textit{(That is why readers can choose the route that matches what they need most. A tired student may begin with Part 1. A teacher reconsidering classroom authority may begin with Part 2. Parents struggling with their children should go directly to Part 3. Senior students or first-year college students may find Part 4 strikes their illusions most directly.)}
\end{truyen}

\section*{Bốn Cửa Đi Vào Quyển Sách}

\begin{tcolorbox}[
  colback=bookbg,
  colframe=bookprimary,
  boxrule=0.8pt,
  arc=4pt,
  title={Lộ Trình Theo Nhu Cầu \quad\textit{(Reading Paths)}},
  fonttitle=\bfseries,
  coltitle=bookprimary,
  breakable,
  enhanced
]
\textbf{Nếu đang thấy mình dễ viện cớ:} đọc Chương 1, 2, 3 rồi quay lại tự hỏi mình đang bảo vệ điều gì.

\textbf{Nếu đang thấy người lớn cũng có phần vô lý:} đọc Chương 4, 5, 6 để học cách phân biệt phản biện tỉnh táo với chống đối cảm tính.

\textbf{Nếu nhà đang căng vì chuyện học:} đọc Chương 7, 8, 9 theo đúng thứ tự, vì ba chương này đi từ ảo tưởng của cha mẹ, đến sự so sánh, rồi tới cuộc va chạm tập thể.

\textbf{Nếu đang bối rối về tương lai:} đọc Chương 10, 11, 12 và ghi lại xem mình đang mê giấc mơ nào nhất nhưng chưa chịu tính chi phí thật.
\end{tcolorbox}

\section*{Đọc Xong Thì Làm Gì?}

\begin{baihoc}
Một quyển sách học đường chỉ thật sự hoàn thành công việc khi nó tạo được một thay đổi nhỏ trong cách người đọc nói chuyện với nhau. Đọc xong một phần, người học nên chọn một câu mình sẽ bớt nói, một thái độ mình sẽ bớt mang vào lớp, và một kỹ năng mình phải tập thêm.
\textit{(A school book truly finishes its work only when it creates a small change in how people speak to each other. After each part, the reader should choose one sentence they will stop saying, one attitude they will bring less into the classroom, and one skill they must practise more.)}
\end{baihoc}

\begin{tuhocnhanh}
1. Nếu chỉ được đọc 3 chương đầu tiên, em sẽ chọn 3 chương nào để chạm đúng vấn đề của mình nhất? \textit{(If you could read only three chapters first, which ones would match your situation best?)}

2. Trong bốn phần của sách, phần nào em nghĩ người lớn nên đọc trước mình? \textit{(Which part should adults read before you do?)}

3. Sau khi đọc xong một phần, em có dám sửa một câu nói quen miệng của mình không? \textit{(After finishing one part, would you dare change one sentence you often say?)}
\end{tuhocnhanh}

\ngancach